\documentclass[10pt]{beamer}
\usetheme{metropolis}
\usepackage{graphicx}\usepackage{listings}\usepackage{booktabs}\usepackage{xcolor}
\definecolor{codebg}{HTML}{F5F5F5}\definecolor{codegreen}{HTML}{2E7D32}\definecolor{codegray}{HTML}{757575}\definecolor{codeblue}{HTML}{1565C0}
\lstdefinestyle{rstyle}{language=R,backgroundcolor=\color{codebg},basicstyle=\ttfamily\scriptsize,keywordstyle=\color{codeblue}\bfseries,commentstyle=\color{codegray}\itshape,breaklines=true,frame=single,rulecolor=\color{codegray!30},showstringspaces=false,morekeywords={library,ggplot,aes,geom_point,geom_hex,geom_density2d,stat_density2d,coord_sf,scale_size,scale_color_viridis,theme_void,labs,sf,st_as_sf,leaflet,addCircleMarkers,addHeatmap}}
\lstdefinestyle{pythonstyle}{language=Python,backgroundcolor=\color{codebg},basicstyle=\ttfamily\scriptsize,keywordstyle=\color{codeblue}\bfseries,commentstyle=\color{codegray}\itshape,breaklines=true,frame=single,rulecolor=\color{codegray!30},showstringspaces=false,morekeywords={import,from,as,gpd,plt,sns,folium,CircleMarker,HeatMap,gaussian_kde,hexbin,scatter}}
\title{Point Maps \& Spatial Patterns}
\subtitle{Workshop 5 --- Module 7}
\author{Prof.Asoc.\ Endri Raco}
\institute{Data Visualization for Data Scientists \\ UNYT Tirana}
\date{2025/26}
\begin{document}
\begin{frame}\titlepage\end{frame}
\begin{frame}{Module Roadmap}
\begin{enumerate}
\item Three point map types: dot, proportional symbol, KDE
\item Overplotting: alpha, hexbin, 2D KDE contour
\item Proportional symbols: encoding a third variable
\item Spatial KDE for hotspot detection
\item Cartograms: distorting geography by data
\end{enumerate}
\begin{exampleblock}{Learning Outcome}
You will produce dot maps for event locations, proportional symbol maps for magnitude, and KDE density surfaces for hotspot detection, handling overplotting with alpha, hexbin, or kernel smoothing.
\end{exampleblock}
\end{frame}
\section{Point Map Types}
\begin{frame}{Dot, Proportional Symbol, KDE}
\begin{center}\includegraphics[width=0.95\textwidth]{figures_w05m07/point_map_types}\end{center}
\end{frame}
\begin{frame}{When to Use Which}
\centering
\begin{tabular}{llp{4.5cm}}
\toprule
\textbf{Type} & \textbf{Encodes} & \textbf{Use When} \\
\midrule
Dot map & Location only & Event mapping (crimes, earthquakes) \\
Proportional symbol & Location + magnitude & City population, store revenue \\
KDE / heatmap & Density of events & Hotspot detection, spatial clustering \\
Hexbin & Binned count per cell & Large $n$ ($>$5K points) \\
\bottomrule
\end{tabular}
\end{frame}
\section{Overplotting}
\begin{frame}{5,000 Points: Three Solutions}
\begin{center}\includegraphics[width=0.95\textwidth]{figures_w05m07/overplotting}\end{center}
\begin{alertblock}{Pro-Tip}
With $>$1,000 points, raw scatter becomes a blob. Three solutions ranked by effectiveness: (1) \textbf{alpha transparency} (cheap, approximate), (2) \textbf{hexbin} (exact counts per cell), (3) \textbf{2D KDE} (smooth continuous density). For maps, KDE is preferred because it produces a continuous surface interpretable as ``event intensity per unit area.''
\end{alertblock}
\end{frame}
\begin{frame}[fragile]{Point Maps in Code}
\begin{columns}[T]
\begin{column}{0.48\textwidth}
\textbf{R}
\begin{lstlisting}[style=rstyle]
# Dot map on sf geometry
ggplot(events_sf) +
  geom_sf(data = districts,
    fill = "grey95") +
  geom_sf(size = 0.5, alpha = 0.3,
    color = "#E53935") +
  theme_void()
# Proportional symbol
ggplot(cities_sf) +
  geom_sf(data = districts) +
  geom_sf(aes(size = population),
    color = "#1565C0", alpha = 0.5) +
  scale_size(range = c(1, 12))
# KDE overlay
ggplot(events_sf) +
  geom_sf(data = districts) +
  stat_density2d(aes(x = lon, y = lat,
    fill = after_stat(level)),
    geom = "polygon", alpha = 0.5) +
  scale_fill_viridis_c()
\end{lstlisting}
\end{column}
\begin{column}{0.48\textwidth}
\textbf{Python}
\begin{lstlisting}[style=pythonstyle]
# Dot map
fig, ax = plt.subplots()
districts.plot(ax=ax, color="grey95",
    edgecolor="grey70")
events.plot(ax=ax, markersize=1,
    color="#E53935", alpha=0.3)
# Proportional symbol
ax.scatter(cities.lon, cities.lat,
    s=cities.pop/1000,
    c="#1565C0", alpha=0.5)
# KDE (scipy)
from scipy.stats import gaussian_kde
xy = np.vstack([events.lon, events.lat])
kde = gaussian_kde(xy)
xg, yg = np.mgrid[...]
z = kde(np.vstack([xg.ravel(),
    yg.ravel()])).reshape(xg.shape)
ax.contourf(xg, yg, z, levels=10,
    cmap="YlOrRd", alpha=0.7)
\end{lstlisting}
\end{column}
\end{columns}
\end{frame}
\section{Cartograms}
\begin{frame}{Cartograms: Geography Distorted by Data}
\centering
\begin{tabular}{lp{4.5cm}l}
\toprule
\textbf{Type} & \textbf{How} & \textbf{R/Python} \\
\midrule
Contiguous & Regions resized but connected & \texttt{cartogram} (R) \\
Dorling & Circles sized by data & \texttt{cartogram::cartogram\_dorling()} \\
Non-contiguous & Shrink/grow maintaining centroid & \texttt{cartogram::cartogram\_ncont()} \\
Tile grid & Each region = equal-size tile & Manual or \texttt{geogrid} \\
\bottomrule
\end{tabular}
\vspace{6pt}
\begin{exampleblock}{Data Insight}
Cartograms solve the choropleth's biggest flaw: large-area, low-population regions dominate visually. By distorting area to match the data variable, cartograms give each region visual weight proportional to its importance.
\end{exampleblock}
\end{frame}
\section{Synthesis}
\begin{frame}{Key Takeaways}
\begin{enumerate}
\item \textbf{Dot maps} for event locations; \textbf{proportional symbols} for location + magnitude; \textbf{KDE} for density/hotspots.
\item \textbf{Overplotting} at $n > 1$K: use alpha, hexbin, or 2D KDE. KDE gives a continuous density surface.
\item \textbf{Cartograms} distort geography so area $\propto$ data, fixing the choropleth's large-region bias.
\item R: \texttt{geom\_sf()} + \texttt{stat\_density2d()} for static, \texttt{leaflet::addHeatmap()} for interactive.
\item Python: \texttt{geopandas.plot()} + \texttt{scipy.stats.gaussian\_kde()} for static, \texttt{folium.plugins.HeatMap()} for interactive.
\end{enumerate}
\end{frame}
\begin{frame}{Essential Readings}
\begin{itemize}
\item Lovelace, R. et al. (2019). \emph{Geocomputation with R}, Chapter 8.
\item Silverman, B.~W. (1986). \emph{Density Estimation for Statistics and Data Analysis}.
\item Dorling, D. (1996). ``Area Cartograms.'' \emph{Concepts and Techniques in Modern Geography}.
\end{itemize}
\vspace{6pt}
\textbf{Next Module}: Interactive Maps (W05-M08)
\end{frame}
\end{document}
